\documentclass[11pt]{article}
\usepackage[margin=1in]{geometry}
\usepackage{graphicx}
\usepackage{amsmath}
\usepackage{pgfplots}
\usepackage{titlesec}
\usepackage{url}
\usepackage{verbatim}
\graphicspath{{../Screenshots/}}
\begin{document}
\section{Course Roadmap}
\subsection{Design Methodology Root cause}
Failure $\rightarrow$ Failure analysis $\rightarrow$ Root cause analysis
\\~\\
Root cause can be traced from these processes:
\begin{itemize}
    \item Design
    \item Material
    \item Manufacturing/Assembly
    \item Service
\end{itemize}
\subsection{5 main failure types}
\begin{enumerate}
    \item distortion    
    \begin{itemize}
        \item yielding
        -Stress state\\
        -Yield criteria (Von mises, Tresca, Rankine)
        \item Buckling
        -$F_{buckling}=\frac{n\pi^2EI}{L^2}$
        \item Creep
        \item Residual stress/relaxation(not covered much in the course)
    \end{itemize}
    \item fracture
    \begin{itemize}
        \item Fractography\\
        -brittle(intragranular - "rock-candy", intergranular - "voids")
        \\-ductile
        \\-fatigue(3 regions: crack initiation, crack propagation(beach marks, striations), brittle fracture)
        \item Stress concentration factor:
        \\
        $k_I=\frac{\sigma_{max}}{\sigma_{app}}=2\sqrt{\frac{a}{\rho}}$
        \item Griffith Criterion
        \\
        $\gamma_s\equiv$ Strain energy
        \\
        $\sigma_c=\sqrt{\frac{2\gamma_sd}{\pi a}}$
        \\
        Compare $k=f\sigma\sqrt{\pi a}$ to $K_{IC}$
        \item fatigue
        \begin{itemize}
            \item $\sigma_{mean}$
            \item S-N curve (Stress life)
            \item 3 Stages (What are these?)
            \item effect of dm?
            \\
            Goodman, Soderberg, Gerber
            \item Paris law
        \end{itemize}
    \end{itemize}
    \item corrosion
    \begin{itemize}
        \item EMF and Galvanic Series
        \\
        -Anodic(gives away electron) vs Cathodic (takes electron)
        \\
        What makes it more anodic?:
        \\
            cold work makes it more anodic?
        \\
            lower oxygen level?
        \item General corrosion
        \item localized corrosion:
        \begin{itemize}
            \item Pitting
            \item Crevice?
            \item Grain boundary
            \item Selective leeching
            \item Stress corrosion cracking
            \\
            Review $\frac{da}{dr}$ vs k plot
        \end{itemize}
        \item velocity affected
        \begin{itemize}
            \item erosion
            \item impingment?
            \item cavitation
        \end{itemize}
        \item Prevention
        \begin{itemize}
            \item databases
            \item testing
            \item design
            \item Cathodic - Using another material to prevent corrosion
            \item anodic - Using another material to prevent corrosion
        \end{itemize}
    \end{itemize}
    \item Wear
    Types of Wear:
    \begin{itemize}
        \item Abrasive
        -Hardness difference between two metals in contact causes plowing/cutting/fracture
        
        \begin{itemize}
            \item 2 body
            \item 3 body\\wear rate can be affected by:\\-particle size $H_a/H_s$\\-microstructure\\environment
        \end{itemize}
        \item Adhesive
        \item Fretting
        \item Lubrication
        \\
        temperature regulation
        \\
        increases gap between surface
        \\
        \textbf{wear surface can be studied using SEM, MEM, etc.}

    \end{itemize}
    \item High T(can be included)
    \begin{itemize}
        \item creep (Study different stages and the creep rate $\epsilon$ vs t plot) insert table from homework on GB climb etc
        \begin{itemize}
            \item Larson Miller Parameter
            \item Failure mechanisms\\
            -grain boundary void diffusion
            \\
            -if climb is dominant, more plastic deformation?
        \end{itemize}
        \item Thermal fatigue
        \item Thermal mechanical fatigue\\-Think of beam in between two walls
        \item Thermal shock
        \\
        $R \propto \frac{\sigma_fk}{\alpha E}$
        \\
        The smaller the thermal expansion the more resistant the material is to thermal shock
        
    \end{itemize}
\end{enumerate}
\subsection{Other}
\begin{enumerate}
    \item NDT
    remember physical aspect on how these methods reveal failure (e.g capillary action)
    \\
    know pros and cons
    \item Manufacturing
\end{enumerate}
\textbf{Focus on welding on the online lecture}
\\
\textbf{There will be a question that relates last question on midterm to fatigue}


\end{document}
